%%%%%%%%%%%%%%%%%%%%%%%%%%%%%%%%%%%%%%%%%%%%%%%%%%%%%%%%%%%%%%%%%%%%%%%%%%%%%%%
% SECTION: THREATS TO VALIDITY (IEEE TSE Format)
%%%%%%%%%%%%%%%%%%%%%%%%%%%%%%%%%%%%%%%%%%%%%%%%%%%%%%%%%%%%%%%%%%%%%%%%%%%%%%%

We discuss threats to the validity of our study following established guidelines
for empirical software engineering research~\cite{arcuri2011practical}.

\subsection{Internal Validity}

\textbf{Implementation Correctness}: We mitigated implementation errors through
unit testing, code review, and comparison with baseline implementations. Our
replication package allows independent verification.

\textbf{Hyperparameter Selection}: Hyperparameters were selected through grid
search on validation data, not test data. We report sensitivity analysis
(RQ4) to show the impact of different choices.

\textbf{Feature Selection}: The reduction from 29 to 19 structural features
was performed on the industrial training set using Random Forest importance
ranking and correlation filtering, without cross-validating the selection
process itself. The Random Forest used only data from the training partition
(temporally preceding the validation and test sets), so the selection did not
observe future build distributions. Nonetheless, this introduces a risk that
the selected subset is over-fitted to the industrial training distribution.
Three factors mitigate this concern: (1)~the sensitivity analysis (RQ4) shows
that the 19-feature set outperforms both the reduced (10) and expanded (29)
sets on temporally held-out validation data; (2)~the same 19 features are
applied to all 20 RTPTorrent projects without re-selection, and the framework
achieves competitive performance across all of them; and (3)~the selected
features are grounded in established TCP domain
knowledge~\cite{chen2023deeporder, kim2002history} (failure rate, recency,
trend, flakiness), reducing the risk that the selection captured
dataset-specific artifacts.

\textbf{Data Leakage}: We ensured strict temporal separation between training
and test data. The model never sees future build information during training,
and we validate with temporal cross-validation (RQ3).

\textbf{Ablation Scope}: We conducted separate ablation studies on both
datasets (Table~\ref{tab:ablation} and Table~\ref{tab:ablation_rtptorrent}).
Because the two evaluation settings use different configurations
(Table~\ref{tab:rtptorrent_config}), the specific components tested differ:
the industrial ablation evaluates the multi-edge graph, orphan handling, and
threshold optimization, while the RTPTorrent ablation evaluates the GATv2,
DNN ensemble, semantic stream, co-failure graph, and DeepOrder features.
The two ablations reveal complementary contribution patterns (graph-dominant
on industrial vs.\ DNN-dominant on RTPTorrent), which we discuss in
Section~\ref{sec:rq2}.

\subsection{External Validity}

\textbf{Dataset Diversity}: We evaluated on two distinct settings, namely an
industrial dataset with 52,102 test executions across 1,339 builds and 20
open-source Java projects from the RTPTorrent benchmark. While this dual
evaluation strengthens generalizability claims, results may not extend to
projects with substantially different testing practices, technology stacks,
or failure distributions.

\textbf{Configuration Customization}: The two evaluation settings required
substantially different configurations of \filopriori{}
(Table~\ref{tab:rtptorrent_config}), including different graph edge types
(5 vs.\ 1), presence or absence of the orphan handling pipeline and DNN
ensemble, and different hyperparameters
(Table~\ref{tab:hyperparams}). While both configurations share the same core
GATv2-based architecture and feature engineering pipeline, deploying
\filopriori{} on a new dataset may require configuration decisions informed by
the available metadata (e.g., whether rich semantic descriptions exist,
whether the test suite evolves frequently enough to produce orphan tests).
We discuss guidelines for these decisions in
Section~\ref{sec:boundary}.

\textbf{Domain Specificity}: The QTA dataset comes from a specific application
domain. Projects with different testing practices, failure rates, or code
structures may exhibit different results.

\textbf{Scale}: Our dataset contains 277 builds with failures. Very large
projects (e.g., Google-scale~\cite{memon2017taming}) may present different
challenges that require additional optimizations.

\textbf{Programming Languages}: All 20 RTPTorrent projects are Java-based,
and the industrial dataset uses a single technology stack. The structural
features (failure rate, recency, execution time) are language-agnostic, but
the semantic stream relies on Sentence-BERT embeddings of test names and
descriptions, whose effectiveness may vary for other languages (e.g.,
Python's descriptive test names vs.\ C++'s terse identifiers) and testing
frameworks (e.g., pytest fixtures, JavaScript async patterns). Furthermore,
the tokenization strategy used for RTPTorrent (camelCase splitting of Java
test names) is language-specific; adapting to other naming conventions
(e.g., Python's snake\_case, C++ namespaces) is straightforward but
untested. Evaluating \filopriori{} on multi-language benchmarks remains
future work.

\subsection{Construct Validity}

\textbf{APFD Metric}: We use APFD as the primary metric, which is standard
in TCP research~\cite{rothermel1999test}. However, APFD assumes equal fault
severity and detection cost. Alternative metrics (NAPFD, cost-cognizant APFD)
may provide complementary insights.

\textbf{Statistical Tests}: We use Wilcoxon signed-rank tests with
$\alpha = 0.05$, which are appropriate for non-parametric comparisons of
paired samples. Following Arcuri and Briand~\cite{arcuri2011practical}, we
report Cliff's delta effect sizes alongside p-values. On the industrial
dataset, effect sizes are negligible to medium (Table~\ref{tab:tcp_comparison}),
indicating that while \filopriori{} consistently outperforms baselines
across builds, the per-build APFD differences are modest. On RTPTorrent,
the limited sample size ($n = 20$ projects) constrains both statistical
power and effect size precision.

\textbf{Baseline Selection}: We compare against five state-of-the-art deep
learning baselines (DeepOrder~\cite{chen2023deeporder},
TCP-Net~\cite{abdelkarim2022tcp}, NodeRank~\cite{li2024noderank},
RETECS~\cite{spieker2017reinforcement}, and
FailRank-BB~\cite{hernandes2025method}), as well as established heuristic
baselines on the RTPTorrent benchmark. While this selection covers diverse
paradigms (pointwise ranking, pairwise ranking, graph-based, reinforcement
learning, and semantic embedding approaches), some recent methods were not
included due to lack of publicly available implementations or incompatible
evaluation settings. Notably, NodeRank~\cite{li2024noderank} was originally
designed for GNN model testing on citation networks (Cora, CiteSeer, PubMed),
not for software TCP. Our adaptation (Section~\ref{sec:related}) applies its
core mutation-based ranking mechanism to the test relationship graph, which is
faithful to the published algorithm but operates in a different domain.
NodeRank's performance in our evaluation may therefore not reflect its
effectiveness in its original setting, and this should be considered when
interpreting comparisons involving NodeRank.

\textbf{Imbalance Handling Across Methods}: Each baseline was trained with its
original class-imbalance strategy rather than a unified protocol. This means
that \filopriori{}'s improvements reflect both architectural advantages and
its specific balancing strategy (balanced sampling + focal loss). We chose
this approach to preserve fidelity to each method's published design, but
acknowledge that a unified balancing protocol could isolate the architectural
contribution more precisely.

To disentangle these contributions, we provide two lines of evidence. First,
the ablation study (Table~\ref{tab:ablation}) shows that removing the
dual-balancing strategy reduces APFD from 0.7611 to 0.7291, which still
exceeds all baselines (the strongest baseline, DeepOrder, achieves 0.6890),
indicating that the architectural advantage persists independently of the
balancing strategy. Second, we conducted a controlled experiment
(Table~\ref{tab:deeporder_dualbalancing}) in which DeepOrder was retrained
using \filopriori{}'s dual-balancing strategy (balanced sampling at 29:1
ratio + focal loss with $\alpha_f = 0.75$, $\gamma = 2.0$), keeping all
other aspects of DeepOrder unchanged. DeepOrder with dual-balancing
achieves APFD = 0.676, \emph{below} its original score (0.689), and 11.3\%
below \filopriori{} (0.761). This confirms that the performance gap is
primarily attributable to the GATv2 graph architecture rather than the
balancing strategy.

\subsection{Conclusion Validity}

\textbf{Statistical Power}: With 277 builds containing failures, we have
sufficient statistical power to detect meaningful differences. Small effect
sizes (e.g., +1.6\% over TCP-Net on RTPTorrent) may not be statistically
significant but can be practically meaningful.

\textbf{Multiple Comparisons}: We compare against five baselines, and
apply Holm-Bonferroni correction for multiple testing in
Section~\ref{sec:rq1_rtptorrent}. On the industrial dataset, all comparisons
yield $p < 0.001$, which survives even conservative Bonferroni correction
($0.001 \times 5 = 0.005 < 0.05$). On RTPTorrent, only NodeRank and
RETECS survive correction; the differences from TCP-Net, FailRank-BB, and
DeepOrder should be interpreted with caution, as detailed in the Results
section.

\subsection{Reproducibility}

To ensure reproducibility, we provide a replication package containing:
\begin{itemize}
    \item Complete source code for \filopriori{} and all five baselines
    \item Configuration files with exact hyperparameters
    (Tables~\ref{tab:hyperparams} and~\ref{tab:arch_params})
    \item Trained model weights for both evaluation settings
    \item Scripts to reproduce all experiments, including ablation,
    temporal cross-validation, and sensitivity analysis
\end{itemize}

\textbf{Dataset availability.} The industrial dataset is subject to a
non-disclosure agreement and cannot be released, either in raw form or as
anonymized derivatives such as pre-computed embeddings. Consequently, the
industrial experiments cannot be reproduced by external researchers, and we
acknowledge that this also limits the community's ability to audit the
semantic stream's behavior on the industrial data. Importantly,
the largest performance margins ($+10.2\%$ to $+27.6\%$ over baselines)
are observed on this proprietary dataset, meaning that the strongest
empirical claims cannot be independently verified from raw data.

To provide tangibility despite the NDA constraints, we include two
representative case studies in Section~\ref{sec:discussion} (a
high-performance build with APFD = 1.0 and a low-performance build with
APFD = 0.152) that describe the model's behavior using masked identifiers
and anonymized feature values. We additionally report aggregate
characterizations of the industrial dataset, including the distribution of
per-build APFD values, structural feature statistics (means, variances,
correlations), and graph connectivity metrics (node degree distributions,
edge type proportions), so that the community can understand the dataset's
properties without access to the raw data.

To compensate for the proprietary dataset limitation, the RTPTorrent
evaluation, which uses a fully public dataset (CC BY 4.0
license~\cite{mattis2020rtptorrent}), serves as the independently
verifiable benchmark. All 20 RTPTorrent projects can be downloaded,
processed, and evaluated end-to-end using the provided scripts without
any access restrictions, allowing full reproduction of every RTPTorrent
result reported in this paper.

